\section{Finiteness conditions}
\label{section:I.6}

\subsection{Noetherian and locally Noetherian preschemes.}
\label{subsection:I.6.1}

\begin{definition}[6.1.1]
\label{I.6.1.1}
We say that a prescheme $X$ is \emph{Noetherian} (resp. \emph{locally
Noetherian}) if it is a finite union (resp. a union) of affine open subsets
$V_\alpha$ such that the ring of each of the schemes induced on the $V_\alpha$
is Noetherian.
\end{definition}

It follows immediately from (\hyperref[I.1.5.2]{1.5.2}) that if $X$ is
locally Noetherian, the structural sheaf $\mathscr{O}_X$ is a
\emph{coherent sheaf of rings}, the question being local. \emph{Every
quasi-coherent sub-$\mathscr{O}_X$-}

\oldpage[I]{141}
\emph{module (resp. every quasi-coherent quotient $\mathscr{O}_X$-module)}
of a \emph{coherent} $\mathscr{O}_X$-module $\mathscr{F}$ is then
\emph{coherent}, because the question is again local, and it suffices
to apply (\hyperref[I.1.5.1]{1.5.1}),
(\hyperref[I.1.4.1]{1.4.1}) and
(\hyperref[I.1.3.10]{1.3.10}), together with the fact that a submodule
(resp. quotient module) of a module of finite type over a Noetherian ring is
of finite type. More particularly, \emph{every quasi-coherent sheaf of
ideals} of $\mathscr{O}_X$ is \emph{coherent}.

If a prescheme $X$ is a finite union (resp. a union) of open subsets
$W_\lambda$ such that the preschemes induced on the $W_\lambda$ are
Noetherian (resp. locally Noetherian), it is clear that $X$ is Noetherian
(resp. locally Noetherian).

\begin{proposition}[6.1.2]
\label{I.6.1.2}
For a prescheme $X$ to be Noetherian, it is necessary and sufficient that it
be locally Noetherian and that its underlying space be quasi-compact.
The underlying space of $X$ is then Noetherian.
\end{proposition}

\begin{proof}
The first assertion follows immediately from the definitions and from
(\hyperref[I.1.1.10]{1.1.10 (ii)}). The second follows from
(\hyperref[I.1.1.6]{1.1.6}) and from the fact that every space that is a
finite union of Noetherian subspaces is Noetherian
(\hyperref[0.2.2.3]{0, 2.2.3}).
\end{proof}

\begin{proposition}[6.1.3]
\label{I.6.1.3}
Let $X$ be an affine scheme with ring $A$. The following conditions are
equivalent: a) $X$ is Noetherian; b) $X$ is locally Noetherian;
c) $A$ is Noetherian.
\end{proposition}

\begin{proof}
The equivalence of a) and b) follows from (\hyperref[I.6.1.2]{6.1.2}) and
from the fact that every affine scheme has a quasi-compact underlying space
(\hyperref[I.1.1.10]{1.1.10}); moreover, it is clear that c) implies a).
To see that a) implies c), note that there is a finite covering
$(V_i)$ of $X$ by affine open subsets such that the ring $A_i$ of the
prescheme induced on $V_i$ is Noetherian. Let then $(\mathfrak{a}_n)$ be an
increasing sequence of ideals of $A$; there corresponds to it canonically and
bijectively (\hyperref[I.1.3.7]{1.3.7}) an increasing sequence
$(\widetilde{\mathfrak{a}}_n)$ of sheaves of ideals in
$\widetilde{A}\equiv\mathscr{O}_X$; to see that the sequence
$(\mathfrak{a}_n)$ is stationary, it suffices to prove that the sequence
$(\widetilde{\mathfrak{a}}_n)$ is. Now the restriction
$\widetilde{\mathfrak{a}}_n|V_i$ is a quasi-coherent sheaf of ideals in
$\mathscr{O}_X|V_i$, being the inverse image of
$\widetilde{\mathfrak{a}}_n$ under the canonical injection $V_i\to X$
(\hyperref[0.5.1.4]{0, 5.1.4}); $\widetilde{\mathfrak{a}}_n|V_i$ is
therefore of the form $\widetilde{\mathfrak{a}}_{ni}$, where $\mathfrak{a}_{ni}$
is an ideal of $A_i$ (\hyperref[I.1.3.7]{1.3.7}). Since $A_i$ is
Noetherian, the sequence $(\mathfrak{a}_{ni})$ is stationary for every $i$,
whence the proposition.
\end{proof}

Note that the preceding argument also proves that \emph{if $X$ is a
Noetherian prescheme, every increasing sequence of coherent sheaves of ideals
of $\mathscr{O}_X$ is stationary}.

\begin{proposition}[6.1.4]
\label{I.6.1.4}
Every sub-prescheme of a Noetherian (resp. locally Noetherian) prescheme
is Noetherian (resp. locally Noetherian).
\end{proposition}

\begin{proof}
It suffices to give the proof for a Noetherian prescheme $X$; moreover,
we are immediately reduced by Definition
(\hyperref[I.6.1.1]{6.1.1}) to the case where $X$ is an affine scheme.
Since every sub-prescheme of $X$ is a closed sub-prescheme of a prescheme
induced on an open subset (\hyperref[I.4.1.3]{4.1.3}), we may restrict to
considering the case of a sub-prescheme $Y$ that is closed or induced on an
open subset of $X$. The case where $Y$ is closed is immediate, because if $A$
is the ring of $X$, we know that $Y$ is an affine scheme with ring
$A/\mathfrak{J}$, where $\mathfrak{J}$ is an ideal of $A$
(\hyperref[I.4.2.3]{4.2.3}); since $A$ is Noetherian
(\hyperref[I.6.1.3]{6.1.3}), so is $A/\mathfrak{J}$.

Suppose now that $Y$ is open in $X$; the underlying space $Y$ is
Noetherian (\hyperref[I.6.1.2]{6.1.2}), hence quasi-compact, and
consequently a finite union of open subsets $D(f_i)$ ($f_i\in A$);
\oldpage[I]{142}
it all comes down to proving the proposition when $Y=D(f)$ with $f\in A$. But
then $Y$ is an affine scheme whose ring is isomorphic to $A_f$
(\hyperref[I.1.3.6]{1.3.6}); since $A$ is Noetherian
(\hyperref[I.6.1.3]{6.1.3}), so is $A_f$.
\end{proof}

\begin{env}[6.1.5]
\label{I.6.1.5}
Note that the \emph{product} of two Noetherian $S$-preschemes is not
necessarily Noetherian, even if these preschemes are affine, because the
tensor product of two Noetherian algebras is not necessarily a Noetherian
ring (cf. (\hyperref[I.6.3.8]{6.3.8})).
\end{env}

\begin{proposition}[6.1.6]
\label{I.6.1.6}
If $X$ is a Noetherian prescheme, the nilradical $\mathscr{N}_X$ of
$\mathscr{O}_X$ is nilpotent.
\end{proposition}

Indeed, we can cover $X$ by finitely many affine open subsets $U_i$, and it
suffices to prove that there exist integers $n_i$ such that
$(\mathscr{N}_X|U_i)^{n_i}=0$; if $n$ is the largest of the $n_i$, we will
then have $\mathscr{N}_X^n=0$. We are therefore reduced to the case where
$X=\operatorname{Spec}(A)$ is affine, $A$ being a Noetherian ring; by
virtue of (\hyperref[I.5.1.1]{5.1.1}) and
(\hyperref[I.1.3.13]{1.3.13}), it suffices to observe that the nilradical
of $A$ is nilpotent (\cite{I-11}, p.~127, cor.~4).

\begin{corollary}[6.1.7]
\label{I.6.1.7}
Let $X$ be a Noetherian prescheme; for $X$ to be an affine scheme, it is
necessary and sufficient that $X_{\mathrm{red}}$ be one.
\end{corollary}

This follows from (\hyperref[I.6.1.6]{6.1.6}) and
(\hyperref[I.5.1.10]{5.1.10}).

\begin{lemma}[6.1.8]
\label{I.6.1.8}
Let $X$ be a topological space, $x$ a point of $X$, and $U$ an open
neighbourhood of $x$ having only finitely many irreducible components. Then
there exists a neighbourhood $V$ of $x$ such that every open neighbourhood of
$x$ contained in $V$ is connected.
\end{lemma}

Let $U_i$ ($1\leq i\leq m$) be the irreducible components of $U$ not
containing $x$; the complement in $U$ of the union of the $U_i$ is an
open neighbourhood $V$ of $x$ in $U$, hence also in $X$; it is moreover
the intersection with $U$ of the complement in $X$ of the union of the
irreducible components of $X$ that do not contain $x$
(\hyperref[0.2.1.6]{0, 2.1.6}). Let then $W$ be an open neighbourhood of
$x$ contained in $V$. The irreducible components of $W$ are the traces on $W$
of the irreducible components of $U$ that meet $W$
(\hyperref[0.2.1.6]{0, 2.1.6}), so these components contain $x$; since
they are connected, so is $W$.

\begin{corollary}[6.1.9]
\label{I.6.1.9}
A locally Noetherian topological space is locally connected (which
implies, among other things, that its connected components are open).
\end{corollary}

\begin{proposition}[6.1.10]
\label{I.6.1.10}
Let $X$ be a locally Noetherian topological space. The following conditions
are equivalent:
\begin{enumerate}
  \item[a)] The irreducible components of $X$ are open.
  \item[b)] The irreducible components of $X$ are identical with its
    connected components.
  \item[c)] The connected components of $X$ are irreducible.
  \item[d)] Two distinct irreducible components of $X$ do not
    meet.
\end{enumerate}

Finally, if $X$ is a prescheme, these conditions are also equivalent to:
\begin{enumerate}
  \item[e)] For every $x\in X$, $\operatorname{Spec}(\mathscr{O}_x)$ is
    irreducible (in other words, the nilradical of $\mathscr{O}_x$ is prime).
\end{enumerate}
\end{proposition}

It is immediate that a) implies b), because an irreducible space is connected,
and a) implies that the irreducible components of $X$ are sets that are both
open and closed. It is trivial that b) implies c); conversely, a
closed set $F$ containing
\oldpage[I]{143}
a connected component $C$ of $X$ and distinct from $C$ cannot be
irreducible, because this set, not being connected, is the union of two
non-empty disjoint sets that are both open and closed in $F$, hence
closed in $X$; consequently c) implies b). We conclude immediately that c)
implies d), two distinct connected components having no point in common.

So far we have not used the fact that $X$ is locally Noetherian.
Suppose now that this hypothesis is satisfied, and let us show that d) implies
a): by virtue of (\hyperref[0.2.1.6]{0, 2.1.6}), we may restrict to the case
where the space $X$ is Noetherian, and so has only finitely many irreducible
components. Since these are closed and pairwise disjoint, they
are open.

Finally, the equivalence of d) and e) holds without assuming that the
underlying space of the prescheme $X$ is locally Noetherian. Indeed, we can
reduce to the case where $X=\operatorname{Spec}(A)$ is affine by virtue of
(\hyperref[0.2.1.6]{0, 2.1.6}); to say that $x$ is contained in only one
irreducible component of $X$ then means that $\mathfrak{j}_x$
contains only one minimal ideal of $A$ (\hyperref[I.1.1.14]{1.1.14}),
which is equivalent to saying that $\mathfrak{j}_x\mathscr{O}_x$ contains only
one minimal ideal of $\mathscr{O}_x$, whence the conclusion.

\begin{corollary}[6.1.11]
\label{I.6.1.11}
Let $X$ be a locally Noetherian space. For $X$ to be irreducible, it is
necessary and sufficient that $X$ be connected and non-empty, and that two
distinct irreducible components of $X$ do not meet. If $X$ is a
prescheme, this last condition is equivalent to
$\operatorname{Spec}(\mathscr{O}_x)$ being irreducible for every $x\in X$.
\end{corollary}

The last part was seen in (\hyperref[I.6.1.10]{6.1.10}); it therefore only
remains to prove the sufficiency of the conditions of the first assertion.
But by (\hyperref[I.6.1.10]{6.1.10}), these conditions imply that
the irreducible components of $X$ are its connected components, and since
$X$ is connected and non-empty, it is irreducible.

\begin{corollary}[6.1.12]
\label{I.6.1.12}
Let $X$ be a locally Noetherian prescheme. For $X$ to be integral, it is
necessary and sufficient that $X$ be connected and non-empty, and that
$\mathscr{O}_x$ be integral for every $x\in X$.
\end{corollary}

\begin{proposition}[6.1.13]
\label{I.6.1.13}
Let $X$ be a locally Noetherian prescheme, and let $x\in X$ be a point such
that the nilradical $\mathscr{N}_x$ of $\mathscr{O}_x$ is prime (resp. such
that $\mathscr{O}_x$ is reduced, resp. integral); then there exists an open
neighbourhood $U$ of $x$ which is irreducible (resp. reduced, resp.
integral).
\end{proposition}

It suffices to consider the two cases where $\mathscr{N}_x$ is prime and where
$\mathscr{N}_x=0$, the third hypothesis being the conjunction of the first
two. If $\mathscr{N}_x$ is prime, $x$ belongs to only one irreducible
component $Y$ of $X$ (\hyperref[I.6.1.10]{6.1.10}); the union of the
irreducible components of $X$ not containing $x$ is closed
(the set of these components being locally finite), and the complement
$U$ of this union is therefore open and contained in $Y$, hence irreducible
(\hyperref[0.2.1.6]{0, 2.1.6}). If $\mathscr{N}_x=0$, we also have
$\mathscr{N}_y=0$ for every $y$ in a neighbourhood of $x$, because
$\mathscr{N}$ is quasi-coherent (\hyperref[I.5.1.1]{5.1.1}), and
consequently coherent since $X$ is locally Noetherian, and the conclusion
follows from (\hyperref[0.5.2.2]{0, 5.2.2}).

\subsection{Artinian preschemes.}
\label{subsection:I.6.2}

\begin{definition}[6.2.1]
\label{I.6.2.1}
We say that a prescheme is \emph{Artinian} if it is affine and its ring is
Artinian.
\end{definition}

\oldpage[I]{144}
\begin{proposition}[6.2.2]
\label{I.6.2.2}
Given a prescheme $X$, the following conditions are equivalent:
\begin{enumerate}
  \item[a)] $X$ is an Artinian scheme;
  \item[b)] $X$ is Noetherian and its underlying space is discrete;
  \item[c)] $X$ is Noetherian and the points of its underlying space are
    closed (condition $\mathrm{T}_1$).
\end{enumerate}

When this is so, the underlying space of $X$ is finite, and the ring $A$
of $X$ is the direct product of the (Artinian) local rings of the points of
$X$.
\end{proposition}

We know that a) implies the last assertion (\cite{I-13}, p.~205, th.~3);
every prime ideal of $A$ is then maximal and is the inverse image of the
maximal ideal of one of the local components of $A$, so the space $X$ is
finite and discrete; a) therefore implies b), and b) evidently implies c). To
see that c) implies a), let us first show that $X$ is then finite; indeed, we
can reduce to the case where $X$ is affine, and we know that a Noetherian
ring all of whose prime ideals are maximal is Artinian (\cite{I-13}, p.~203),
whence our assertion. The underlying space $X$ is then discrete, the
topological sum of finitely many points $x_i$, and the local rings
$\mathscr{O}_{x_i}=A_i$ are Artinian; it is clear that $X$ is isomorphic to
the affine scheme that is the prime spectrum of the ring $A$, the direct
product of the $A_i$ (\hyperref[I.1.7.3]{1.7.3}).

\subsection{Morphisms of finite type.}
\label{subsection:I.6.3}

\begin{definition}[6.3.1]
\label{I.6.3.1}
We say that a morphism $f:X\to Y$ is \emph{of finite type} if $Y$ is the union
of a family $(V_\alpha)$ of affine open subsets having the following property:
\begin{enumerate}
  \item[(P)] $f^{-1}(V_\alpha)$ is a finite union of affine open subsets
    $U_{\alpha i}$ such that each of the rings $A(U_{\alpha i})$ is an
    algebra of finite type over $A(V_\alpha)$.
\end{enumerate}
We then also say that $X$ is a prescheme of finite type over $Y$, or a
$Y$-prescheme of finite type.
\end{definition}

\begin{proposition}[6.3.2]
\label{I.6.3.2}
If $f:X\to Y$ is a morphism of finite type, every affine open subset $W$ of
$Y$ has the property (P) of Definition (\hyperref[I.6.3.1]{6.3.1}).
\end{proposition}

We will first prove the

\begin{lemma}[6.3.2.1]
\label{I.6.3.2.1}
If $T\subset Y$ is an affine open subset having the property (P), then, for
every $g\in A(T)$, $D(g)$ has the property (P).
\end{lemma}

Indeed, by hypothesis, $f^{-1}(T)$ is a finite union of affine open subsets
$Z_j$ such that $A(Z_j)$ is an algebra of finite type over $A(T)$; let
$\varphi_j:A(T)\to A(Z_j)$ be the ring homomorphism corresponding to the
restriction of $f$ to $Z_j$ (\hyperref[I.2.2.4]{2.2.4}), and set
$g_j=\varphi_j(g)$; we then have
$f^{-1}(D(g))\cap Z_j=D(g_j)$
(\hyperref[I.1.2.2.2]{1.2.2.2}). Now
$A(D(g_j))=A(Z_j)_{g_j}=A(Z_j)[1/g_j]$ is of finite type over $A(Z_j)$, and
\emph{a fortiori} over $A(T)$ by virtue of the hypothesis, hence also over
$A(D(g))=A(T)[1/g]$, which proves the lemma.

This lemma being established, since $W$ is quasi-compact
(\hyperref[I.1.1.10]{1.1.10}), there exists a finite covering of $W$ by
sets of the form $D(g_i)$, where each $g_i$ belongs to a ring
$A(V_{\alpha(i)})$. Each $D(g_i)$, being quasi-compact, is a finite union
of sets $D(h_{ik})$ where $h_{ik}\in A(W)$; if
$\varphi_i:A(W)\to A(D(g_i))$ is the canonical map, we have
$D(h_{ik})=D(\varphi_i(h_{ik}))$ by virtue of
(\hyperref[I.1.2.2.2]{1.2.2.2}). By virtue of
(\hyperref[I.6.3.2.1]{6.3.2.1}), each of the
$f^{-1}(D(h_{ik}))$ admits a finite covering by affine open subsets
$U_{ijk}$ such that $A(U_{ijk})$ is an algebra of finite type over
$A(D(h_{ik}))=A(W)[1/h_{ik}]$, whence the proposition.

\oldpage[I]{145}
We can therefore say that the notion of prescheme of finite type over $Y$ is
\emph{local on $Y$}.

\begin{proposition}[6.3.3]
\label{I.6.3.3}
Let $X$, $Y$ be two affine schemes; for $X$ to be of finite type over
$Y$, it is necessary and sufficient that $A(X)$ be an algebra of finite type
over $A(Y)$.
\end{proposition}

The condition being evidently sufficient, let us prove that it is necessary.
Set $A=A(Y)$, $B=A(X)$; by virtue of
(\hyperref[I.6.3.2]{6.3.2}), there exists a finite affine open covering
$(V_i)$ of $X$ such that each of the rings $A(V_i)$ is an $A$-algebra of
finite type. Moreover, the $V_i$ being quasi-compact, we can cover each of
them by finitely many open subsets of the form $D(g_{ij})\subset V_i$, where
$g_{ij}\in B$; if $\varphi_i$ is the homomorphism $B\to A(V_i)$ corresponding
to the canonical injection $V_i\to X$, we have
\[
  B_{g_{ij}}=(A(V_i))_{\varphi_i(g_{ij})}
  =A(V_i)[1/\varphi_i(g_{ij})],
\]
so $B_{g_{ij}}$ is an $A$-algebra of finite type. We can therefore reduce
to the case where $V_i=D(g_i)$ with $g_i\in B$. By hypothesis, there exist a
finite subset $F_i$ of $B$ and an integer $n_i\geq 0$ such that $B_{g_i}$ is
the algebra generated over $A$ by the elements $b_i/g_i^{n_i}$, where $b_i$
runs through $F_i$. Since there are only finitely many $g_i$, we may moreover
suppose all the $n_i$ equal to one and the same integer $n$. Furthermore,
since the $D(g_i)$ form a covering of $X$, the ideal generated in $B$ by the
$g_i$ is equal to $B$; in other words, there exist $h_i\in B$ such that
$\sum_i h_i g_i=1$. Let then $F$ be the finite subset of $B$ that is the union
of the $F_i$, of the set of the $g_i$ and of the set of the $h_i$; let us show
that the subring $B'=A[F]$ of $B$ is equal to $B$. By hypothesis, for every
$b\in B$ and every $i$, the canonical image of $b$ in $B_{g_i}$ is of the form
$b'_i/g_i^{m_i}$, where $b'_i\in B'$; by multiplying the $b'_i$ by suitable
powers of the $g_i$, we may further suppose all the $m_i$ equal to one and
the same integer $m$. By the definition of rings of fractions, there is
therefore an integer $N$ (depending on $b$) such that $N\geq m$ and
$g_i^N b\in B'$ for every $i$; now, \emph{in the ring $B'$}, the $g_i^N$
generate the ideal $B'$, since the $g_i$ do (the $h_i$ belonging to $B'$);
there are therefore $c_i\in B'$ such that $\sum_i c_i g_i^N=1$, whence
$b=\sum_i c_i g_i^N b\in B'$, Q.E.D.

\begin{proposition}[6.3.4]
\label{I.6.3.4}
\begin{enumerate}
  \item[(i)] Every closed immersion is of finite type.
  \item[(ii)] The composite of two morphisms of finite type is of finite
    type.
  \item[(iii)] If $f:X\to X'$, $g:Y\to Y'$ are two $S$-morphisms of finite
    type, $f\times_S g$ is of finite type.
  \item[(iv)] If $f:X\to Y$ is an $S$-morphism of finite type, $f_{(S')}$ is
    of finite type for every extension $g:S'\to S$ of the base prescheme.
  \item[(v)] If the composite $g\circ f$ of two morphisms is of finite type,
    and if $g$ is separated, $f$ is of finite type.
  \item[(vi)] If a morphism $f$ is of finite type, then so is
    $f_{\mathrm{red}}$.
\end{enumerate}
\end{proposition}

By virtue of (\hyperref[I.5.5.12]{5.5.12}), it suffices to prove (i),
(ii) and (iv).

To establish (i), we may restrict to the case of a canonical injection
$X\to Y$, $X$ being a closed sub-prescheme of $Y$; moreover
(\hyperref[I.6.3.2]{6.3.2}), we may suppose $Y$ affine, in which case $X$
is also affine (\hyperref[I.4.2.3]{4.2.3}) and its ring is isomorphic to
a quotient ring $A/\mathfrak{J}$, where $A$ is the ring of $Y$ and
$\mathfrak{J}$ an ideal of $A$; since $A/\mathfrak{J}$ is of finite type over
$A$, the conclusion follows.

Let us now prove (ii). Let $f:X\to Y$, $g:Y\to Z$ be two morphisms of
finite type, and let $U$ be an affine open subset of $Z$; $g^{-1}(U)$ admits a
finite covering by affine open subsets $V_i$ such that $A(V_i)$ is an
algebra of finite type over $A(U)$ (\hyperref[I.6.3.2]{6.3.2}); similarly,

\oldpage[I]{146}
each of the $f^{-1}(V_i)$ admits a finite covering by affine open subsets
$W_{ij}$ such that $A(W_{ij})$ is an algebra of finite type over $A(V_i)$, and
consequently also an algebra of finite type over $A(U)$; whence the
conclusion.

Finally, to prove (iv), we may restrict to the case where $S=Y$; indeed,
$f_{(S')}$ is also equal to $f_{(Y_{(S')})}$, $f$ being considered as a
$Y$-morphism and the base extension being $Y_{(S')}\to Y$
(\hyperref[I.3.3.9]{3.3.9}). Let then $p$, $q$ be the projections
$X_{(S')}\to X$ and $X_{(S')}\to S'$. Let $V$ be an affine open subset in
$S$; $f^{-1}(V)$ is a finite union of affine open subsets $W_i$ each of which
is such that $A(W_i)$ is an algebra of finite type over $A(V)$
(\hyperref[I.6.3.2]{6.3.2}). Let $V'$ be an affine open subset of $S'$
contained in $g^{-1}(V)$; since $f\circ p=g\circ q$, $q^{-1}(V')$ is contained
in the union of the $p^{-1}(W_i)$; on the other hand, the intersection
$p^{-1}(W_i)\cap q^{-1}(V')$ is identified with the product $W_i\times_V V'$
(\hyperref[I.3.2.7]{3.2.7}), which is an affine scheme with ring
isomorphic to $A(W_i)\otimes_{A(V)}A(V')$ (\hyperref[I.3.2.2]{3.2.2});
since the latter is by hypothesis an algebra of finite type over $A(V')$, the
proposition is proved.

\begin{corollary}[6.3.5]
\label{I.6.3.5}
Let $f:X\to Y$ be an immersion morphism. If the underlying space of $Y$
(resp. $X$) is locally Noetherian (resp. Noetherian), $f$ is of finite type.
\end{corollary}

We may always suppose $Y$ affine (\hyperref[I.6.3.2]{6.3.2}); if
the underlying space of $Y$ is locally Noetherian, we may moreover
suppose it Noetherian, and then the underlying space of $X$, which is a
subspace of it, is Noetherian. In other words, we may suppose $Y$ affine and
the underlying space of $X$ Noetherian; there then exists a covering of $X$
by finitely many affine open subsets $D(g_i)\subset Y$, where $g_i\in A(Y)$,
such that $X\cap D(g_i)$ is closed in $D(g_i)$ (hence an affine scheme
(\hyperref[I.4.2.3]{4.2.3})), since $X$ is locally closed in $Y$
(\hyperref[I.4.1.3]{4.1.3}). Then $A(X\cap D(g_i))$ is an algebra of
finite type over $A(D(g_i))$, by (\hyperref[I.6.3.4]{6.3.4, (i)}) and
(\hyperref[I.6.3.3]{6.3.3}), and
$A(D(g_i))=A(Y)_{g_i}=A(Y)[1/g_i]$ is of finite type over $A(Y)$, which
completes the proof.

\begin{corollary}[6.3.6]
\label{I.6.3.6}
Let $f:X\to Y$, $g:Y\to Z$ be two morphisms. If $g\circ f$ is of finite type,
and if $X$ is Noetherian, or $X\times_Z Y$ locally Noetherian, then $f$ is of
finite type.
\end{corollary}

This follows immediately from the proof of
(\hyperref[I.5.5.12]{5.5.12}) and from (\hyperref[I.6.3.5]{6.3.5})
applied to the immersion morphism $\Gamma_f$.

\begin{proposition}[6.3.7]
\label{I.6.3.7}
Let $f:X\to Y$ be a morphism of finite type; if $Y$ is Noetherian (resp.
locally Noetherian), then $X$ is Noetherian (resp. locally Noetherian).
\end{proposition}

We may restrict to giving the proof when $Y$ is Noetherian. Then $Y$
is a finite union of affine open subsets $V_i$ such that the $A(V_i)$ are
Noetherian rings. By virtue of (\hyperref[I.6.3.2]{6.3.2}), each of the
$f^{-1}(V_i)$ is a union of finitely many affine open subsets $W_{ij}$ such
that the $A(W_{ij})$ are algebras of finite type over $A(V_i)$, hence
Noetherian rings; this proves that $X$ is Noetherian.

\begin{corollary}[6.3.8]
\label{I.6.3.8}
Let $X$ be a prescheme of finite type over $S$. For every base extension
$S'\to S$ such that $S'$ is Noetherian (resp. locally Noetherian),
$X_{(S')}$ is Noetherian (resp. locally Noetherian).
\end{corollary}

This follows from (\hyperref[I.6.3.7]{6.3.7}), $X_{(S')}$ being of finite
type over $S'$ by virtue of (\hyperref[I.6.3.4]{6.3.4, (iv)}).

We can also say that in a product $X\times_S Y$ of $S$-preschemes, if
\emph{one} of the
\oldpage[I]{147}
factors $X$, $Y$ is \emph{of finite type} over $S$ and
\emph{the other Noetherian} (resp. \emph{locally Noetherian}), then
$X\times_S Y$ is \emph{Noetherian} (resp. \emph{locally Noetherian}).

\begin{corollary}[6.3.9]
\label{I.6.3.9}
Let $X$ be a prescheme of finite type over a locally Noetherian prescheme
$S$. Then every $S$-morphism $f:X\to Y$ is of finite type.
\end{corollary}

Indeed, we may suppose $S$ Noetherian; if $\varphi:X\to S$,
$\psi:Y\to S$ are the structural morphisms, we have $\varphi=\psi\circ f$, and
$X$ is Noetherian by virtue of (\hyperref[I.6.3.7]{6.3.7}); $f$ is
therefore of finite type by virtue of (\hyperref[I.6.3.6]{6.3.6}).

\begin{proposition}[6.3.10]
\label{I.6.3.10}
Let $f:X\to Y$ be a morphism of finite type. For $f$ to be surjective, it is
necessary and sufficient that, for every algebraically closed field $\Omega$,
the map $X(\Omega)\to Y(\Omega)$ corresponding to $f$
(\hyperref[I.3.4.1]{3.4.1}) be surjective.
\end{proposition}

The condition is sufficient, as we see by considering, for every
$y\in Y$, an algebraically closed extension $\Omega$ of $k(y)$, and the
commutative diagram
\[
  \xymatrix{
    X\ar[dd]_f\\
    & \operatorname{Spec}(\Omega)\ar[ul]\ar[dl]\\
    Y
  }
\]
(\emph{cf.} (\hyperref[I.3.5.3]{3.5.3})). Conversely, suppose $f$
surjective, and let $g:\{\xi\}=\operatorname{Spec}(\Omega)\to Y$ be a
morphism, $\Omega$ being an algebraically closed field. If we consider the
diagram
\[
  \xymatrix{
    X\ar[d]_f &
    X_{(\Omega)}\ar[l]\ar[d]^{f_{(\Omega)}}\\
    Y &
    \operatorname{Spec}(\Omega)\ar[l]
  }
\]
it therefore suffices to show that there exists in $X_{(\Omega)}$ a
\emph{point rational over $\Omega$}
(\hyperref[I.3.3.14]{3.3.14}, \hyperref[I.3.4.3]{3.4.3} and
\hyperref[I.3.4.4]{3.4.4}). Since $f$ is surjective, $X_{(\Omega)}$ is
not empty (\hyperref[I.3.5.10]{3.5.10}), and since $f$ is of finite type,
so is $f_{(\Omega)}$ (\hyperref[I.6.3.4]{6.3.4, (iv)}); hence
$X_{(\Omega)}$ contains a non-empty affine open subset $Z$ such that $A(Z)$ is
a nonzero algebra of finite type over $\Omega$. By virtue of Hilbert's
Nullstellensatz \cite{I-21}, there exists an $\Omega$-homomorphism
$A(Z)\to\Omega$, hence a section of $X_{(\Omega)}$ over
$\operatorname{Spec}(\Omega)$, which proves the proposition.

\subsection{Algebraic preschemes.}
\label{subsection:I.6.4}

\begin{definition}[6.4.1]
\label{I.6.4.1}
Given a field $K$, we call a prescheme $X$ of finite type over $K$ an
\emph{algebraic $K$-prescheme}; $K$ is called the \emph{base field}
of $X$. If moreover $X$ is a scheme (or, what amounts to the same thing
(\hyperref[I.5.5.8]{5.5.8}), if $X$ is a $K$-scheme), we also say that
$X$ is an \emph{algebraic $K$-scheme}.
\end{definition}

Every algebraic $K$-prescheme is \emph{Noetherian}
(\hyperref[I.6.3.7]{6.3.7}).

\begin{proposition}[6.4.2]
\label{I.6.4.2}
Let $X$ be an algebraic $K$-prescheme. For a point $x\in X$ to be closed, it
is necessary and sufficient that $k(x)$ be an algebraic extension of $K$ of
finite degree.
\end{proposition}

We may suppose $X$ affine, the ring $A$ of $X$ being a $K$-algebra of finite
type. Indeed, the affine open subsets $U$ of $X$ such that $A(U)$ is a
$K$-algebra of finite type form a covering of $X$
(\hyperref[I.6.3.1]{6.3.1}). The closed points of $X$ are then the points
such that $\mathfrak{j}_x$ is
\oldpage[I]{148}
a maximal ideal of $A$, in other words such that $A/\mathfrak{j}_x$ is a
field (necessarily equal to $k(x)$). Since $A/\mathfrak{j}_x$ is a
$K$-algebra of finite type, we see that if $x$ is closed, $k(x)$ is a field
which is an algebra of finite type over $K$, hence necessarily a $K$-algebra
of \emph{finite rank} \cite{I-21}. Conversely, if $k(x)$ is of finite rank
over $K$, so is $A/\mathfrak{j}_x\subset k(x)$, and since every integral ring
that is a $K$-algebra of finite rank is a field, we have
$A/\mathfrak{j}_x=k(x)$, so $x$ is closed.

\begin{corollary}[6.4.3]
\label{I.6.4.3}
Let $K$ be an algebraically closed field and $X$ an algebraic $K$-prescheme;
the closed points of $X$ are then the points rational over $K$
(\hyperref[I.3.4.4]{3.4.4}) and are canonically identified with the points
of $X$ with values in $K$.
\end{corollary}

\begin{proposition}[6.4.4]
\label{I.6.4.4}
Let $X$ be an algebraic prescheme over a field $K$. The following properties
are equivalent:
\begin{enumerate}
  \item[\rm(a)] $X$ is Artinian.
  \item[\rm(b)] The underlying space of $X$ is discrete.
  \item[\rm(c)] The underlying space of $X$ has only finitely many closed
    points.
  \item[\rm(c')] The underlying space of $X$ is finite.
  \item[\rm(d)] The points of $X$ are closed.
  \item[\rm(e)] $X$ is isomorphic to $\operatorname{Spec}(A)$, where $A$ is a
    $K$-algebra of finite rank.
\end{enumerate}
\end{proposition}

Since $X$ is Noetherian, it follows from (\hyperref[I.6.2.2]{6.2.2}) that
the conditions a), b), d) are equivalent and imply c) and c'); moreover,
it is clear that e) implies a). It remains to see that c) implies d) and
e); we may restrict to the case where $X$ is affine. Then $A(X)$ is a
$K$-algebra of finite type (\hyperref[I.6.3.3]{6.3.3}), hence a Jacobson
ring (\cite[p.~3-11 and 3-12]{I-1}), in which there are by hypothesis
only finitely many maximal ideals. Since a finite intersection of prime
ideals can be a prime ideal only if it is equal to one of them,
every prime ideal of $A(X)$ is therefore maximal, whence d). Moreover, we
then know (\hyperref[I.6.2.2]{6.2.2}) that $A(X)$ is an Artinian
$K$-algebra of finite type, hence necessarily of \emph{finite rank}
\cite{I-21}.

\begin{env}[6.4.5]
\label{I.6.4.5}
When the conditions of (\hyperref[I.6.4.4]{6.4.4}) are satisfied,
we say that $X$ is a scheme \emph{finite over $K$}
(\emph{cf.} (\hyperref[II.6.1.1]{II, 6.1.1})), or a
\emph{finite $K$-scheme}, of \emph{rank} $[A:K]$, which we also denote by
$\operatorname{rg}_K(X)$; if $X$, $Y$ are two schemes finite over $K$, we
have
\[
  \operatorname{rg}_K(X\amalg Y)
  =\operatorname{rg}_K(X)+\operatorname{rg}_K(Y)
  \tag{6.4.5.1}\label{I.6.4.5.1}
\]
\[
  \operatorname{rg}_K(X\times_K Y)
  =\operatorname{rg}_K(X)\operatorname{rg}_K(Y)
  \tag{6.4.5.2}\label{I.6.4.5.2}
\]
as follows from (\hyperref[I.3.2.2]{3.2.2}).
\end{env}

\begin{corollary}[6.4.6]
\label{I.6.4.6}
Let $X$ be a scheme finite over a field $K$. For every extension $K'$ of $K$,
$X\otimes_K K'$ is a scheme finite over $K'$, and its rank over $K'$ is equal
to the rank of $X$ over $K$.
\end{corollary}

Indeed, if $A=A(X)$, we have $[A\otimes_K K':K']=[A:K]$.

\begin{corollary}[6.4.7]
\label{I.6.4.7}
Let $X$ be a scheme finite over a field $K$; set
$n=\sum_{x\in X}[k(x):K]_s$
\emph{(recall that if $K'$ is an extension of $K$, $[K':K]_s$ is the
\emph{separable rank} of $K'$ over $K$, the rank of the largest separable
algebraic extension of $K$ contained in $K'$)};
\oldpage[I]{149}
then, for every algebraically closed extension $\Omega$ of $K$, the
underlying space of $X\otimes_K\Omega$ has exactly $n$ points, which are
identified with the points of $X$ with values in $\Omega$.
\end{corollary}

We may evidently restrict to the case where the ring $A=A(X)$ is local
(\hyperref[I.6.2.2]{6.2.2}); let $\mathfrak m$ be its maximal ideal and
$L=A/\mathfrak m$ its residue field, an algebraic extension of $K$. The points
of $X$ with values in $\Omega$ then correspond bijectively to the
$\Omega$-sections of $X\otimes_K\Omega$
(\hyperref[I.3.4.1]{3.4.1} and \hyperref[I.3.3.14]{3.3.14}), and also
to the $K$-homomorphisms of $L$ into $\Omega$
(\hyperref[I.1.7.3]{1.7.3}), whence the proposition (Bourbaki, \emph{Alg.},
chap.~V, \textsection7, n\textsuperscript{o}~5, prop.~8), taking into
account (\hyperref[I.6.4.3]{6.4.3}).

\begin{env}[6.4.8]
\label{I.6.4.8}
The number $n$ defined in (\hyperref[I.6.4.7]{6.4.7}) is called the
\emph{separable rank} of $A$ (or of $X$) over $K$, or also the
\emph{geometric number of points of $X$}; it is therefore equal to the number
of elements of $X(\Omega)_K$. It follows immediately from this definition
that, for every extension $K'$ of $K$, $X\otimes_K K'$ has the same geometric
number of points as $X$. If we denote this number by $n(X)$, it is clear that
if $X$, $Y$ are two schemes finite over $K$, we have
\[
  n(X\amalg Y)=n(X)+n(Y).
  \tag{6.4.8.1}\label{I.6.4.8.1}
\]

Under the same hypotheses, we also have
\[
  n(X\times_K Y)=n(X)n(Y)
  \tag{6.4.8.2}\label{I.6.4.8.2}
\]
as follows immediately from the interpretation of $n(X)$ as the number
of elements of $X(\Omega)_K$ and from formula
(\hyperref[I.3.4.3.1]{3.4.3.1}).
\end{env}

\begin{proposition}[6.4.9]
\label{I.6.4.9}
Let $K$ be a field, $X$, $Y$ two algebraic $K$-preschemes, $f:X\to Y$ a
$K$-morphism, and $\Omega$ an algebraically closed extension of $K$ of
infinite transcendence degree over $K$. For $f$ to be surjective, it is
necessary and sufficient that the map $X(\Omega)_K\to Y(\Omega)_K$
corresponding to $f$ (\hyperref[I.3.4.1]{3.4.1}) be surjective.
\end{proposition}

The necessity follows from (\hyperref[I.6.3.10]{6.3.10}), noting that
$f$ is necessarily of finite type (\hyperref[I.6.3.9]{6.3.9}). To see
that the condition is sufficient, we argue as in
(\hyperref[I.6.3.10]{6.3.10}), noting that for every $y\in Y$,
$k(y)$ is an extension of $K$ of finite type, and consequently is
$K$-isomorphic to a subfield of $\Omega$.

\begin{remark}[6.4.10]
\label{I.6.4.10}
We will see in Chapter~IV that the conclusion of
(\hyperref[I.6.4.9]{6.4.9}) still holds without any hypothesis on the
transcendence degree of $\Omega$ over $K$.
\end{remark}

\begin{proposition}[6.4.11]
\label{I.6.4.11}
If $f:X\to Y$ is a morphism of finite type, then for every $y\in Y$, the fibre
$f^{-1}(y)$ is an algebraic prescheme over the residue field $k(y)$, and for
every $x\in f^{-1}(y)$, $k(x)$ is an extension of finite type of $k(y)$.
\end{proposition}

Since $f^{-1}(y)=X\otimes_Y k(y)$
(\hyperref[I.3.6.3]{3.6.3}), the proposition follows from
(\hyperref[I.6.3.4]{6.3.4, (iv)}) and
(\hyperref[I.6.3.3]{6.3.3}).

\begin{proposition}[6.4.12]
\label{I.6.4.12}
Let $f:X\to Y$, $g:Y'\to Y$ be two morphisms; set
$X'=X\times_Y Y'$ and let $f'=f_{(Y')}:X'\to Y'$. Let $y'\in Y'$ and
$y=g(y')$; if the fibre $f^{-1}(y)$ is a finite algebraic scheme over $k(y)$,
then the fibre $f'^{-1}(y')$ is a finite algebraic scheme over $k(y')$, having
the same rank and the same geometric number of points as $f^{-1}(y)$.
\end{proposition}

Taking into account the transitivity of fibres
(\hyperref[I.3.6.5]{3.6.5}), this follows immediately from
(\hyperref[I.6.4.6]{6.4.6}) and
(\hyperref[I.6.4.8]{6.4.8}).

\begin{env}[6.4.13]
\label{I.6.4.13}
Proposition (\hyperref[I.6.4.11]{6.4.11}) shows that morphisms of finite
type correspond intuitively to ``algebraic families of algebraic
varieties'', the points of $Y$ playing
\oldpage[I]{150}
the role of ``parameters'', which gives these morphisms a
``geometric'' meaning. Morphisms that are not of finite type will occur
later mainly in questions of ``change of the base prescheme'', for
example by localization or completion.
\end{env}


\subsection{Local determination of a morphism}
\label{subsection:I.6.5}

\begin{proposition}[6.5.1]
\label{I.6.5.1}
Let $X$, $Y$ be two $S$-preschemes, $Y$ being of finite type over $S$; let
$x\in X$, $y\in Y$ lie above the same point $s\in S$.
\begin{enumerate}
  \item[(i)] If two $S$-morphisms $f=(\psi,\theta)$,
    $f'=(\psi',\theta')$ from $X$ to $Y$ are such that
    $\psi(x)=\psi'(x)=y$ and the (local) $\mathscr{O}_s$-homomorphisms
    $\theta_x^\sharp$ and ${\theta'}_x^\sharp$ from $\mathscr{O}_y$ to
    $\mathscr{O}_x$ are identical, then $f$ and $f'$ coincide in an open
    neighbourhood of $x$.
  \item[(ii)] Suppose moreover that $S$ is locally Noetherian. For every local
    $\mathscr{O}_s$-homomorphism
    $\varphi:\mathscr{O}_y\to\mathscr{O}_x$, there exist an open neighbourhood
    $U$ of $x$ in $X$ and an $S$-morphism $f=(\psi,\theta)$ from $U$ to $Y$
    such that $\psi(x)=y$ and $\theta_x^\sharp=\varphi$.
\end{enumerate}
\end{proposition}

\begin{enumerate}
  \item[(i)] The question being local on $S$, $X$ and $Y$, we may suppose
    $S$, $X$, $Y$ affine with respective rings $A$, $B$, $C$, $f$ and $f'$
    being of the form $({}^a\varphi,\widetilde{\varphi})$ and
    $({}^a\varphi',\widetilde{\varphi}')$ respectively, where $\varphi$ and
    $\varphi'$ are two $A$-homomorphisms from $C$ to $B$ such that
    $\varphi^{-1}(\mathfrak{j}_x)={\varphi'}^{-1}(\mathfrak{j}_x)
    =\mathfrak{j}_y$, and the homomorphisms $\varphi_x$ and $\varphi'_x$ from
    $C_y$ to $B_x$ deduced from $\varphi$ and $\varphi'$ are identical;
    we may moreover suppose that $C$ is an $A$-algebra of finite type. Let
    $c_i$ ($1\leq i\leq n$) be generators of the $A$-algebra $C$, and set
    $b_i=\varphi(c_i)$, $b'_i=\varphi'(c_i)$; by hypothesis, we have
    $b_i/1=b'_i/1$ in the ring of fractions $B_x$ ($1\leq i\leq n$). This
    means that there exist elements $s_i\in B-\mathfrak{j}_x$ such that
    $s_i(b_i-b'_i)=0$ for $1\leq i\leq n$, and we may evidently suppose
    all the $s_i$ equal to one and the same element $g\in B-\mathfrak{j}_x$. We
    conclude that $b_i/1=b'_i/1$ for $1\leq i\leq n$ in the ring of
    fractions $B_g$; if $i_g$ is the canonical homomorphism $B\to B_g$, we
    consequently have $i_g\circ\varphi=i_g\circ\varphi'$; hence the
    restrictions of $f$ and $f'$ to $D(g)$ are identical.
  \item[(ii)] We may reduce to the same situation as in (i), and suppose
    moreover that the ring $A$ is Noetherian. Let $c_i$
    ($1\leq i\leq n$) be generators of the $A$-algebra $C$, and let
    $\alpha:A[X_1,\ldots,X_n]\to C$ be the homomorphism of the polynomial
    algebra $A[X_1,\ldots,X_n]$ onto $C$ sending $X_i$ to $c_i$ for
    $1\leq i\leq n$. On the other hand, let $i_y$ be the canonical homomorphism
    $C\to C_y$, and consider the composite homomorphism
    \[
      \beta:A[X_1,\ldots,X_n]
      \xrightarrow{\alpha} C
      \xrightarrow{i_y} C_y
      \xrightarrow{\varphi} B_x.
    \]
    Denote by $\mathfrak{a}$ the kernel of $\beta$; since $A$ is
    Noetherian, so is $A[X_1,\ldots,X_n]$, and consequently
    $\mathfrak{a}$ admits a finite system of generators
    $Q_j(X_1,\ldots,X_n)$ ($1\leq j\leq m$). On the other hand, each of the
    elements $\varphi(i_y(c_i))$ can be written $b_i/s_i$, where $b_i\in B$ and
    $s_i\notin\mathfrak{j}_x$; we may moreover suppose all the $s_i$ equal
    to one and the same element $g\in B-\mathfrak{j}_x$. This being so, we have
    by hypothesis $Q_j(b_1/g,\ldots,b_n/g)=0$ in $B_x$; set
    \[
      Q_j(X_1/T,\ldots,X_n/T)
      =P_j(X_1,\ldots,X_n,T)/T^{k_j}
    \]
    where $P_j$ is homogeneous of degree $k_j$. Then let
    $d_j=P_j(b_1,\ldots,b_n,g)\in B$. By hypothesis, we have $t_jd_j=0$ for some
    $t_j\in B-\mathfrak{j}_x$ ($1\leq j\leq m$), and we may evidently
    suppose all the $t_j$
\oldpage[I]{151}
    equal to one and the same element $h\in B-\mathfrak{j}_x$; we conclude that
    $P_j(hb_1,\ldots,hb_n,hg)=0$ for $1\leq j\leq m$. This being so, consider
    the homomorphism $\rho$ of $A[X_1,\ldots,X_n]$ into the ring of fractions
    $B_{hg}$ which maps $X_i$ to $hb_i/hg$ ($1\leq i\leq n$); the image of
    $\mathfrak{a}$ under this homomorphism is 0, and \emph{a fortiori} so is
    the image under $\rho$ of the kernel $\alpha^{-1}(0)$. Hence $\rho$
    factors as
    $A[X_1,\ldots,X_n]\xrightarrow{\alpha}C\xrightarrow{\gamma}B_{hg}$, with
    $\gamma(c_i)=hb_i/hg$, and it is clear that if $i_x$ is the canonical
    homomorphism $B_{hg}\to B_x$, the diagram
    \[
      \xymatrix{
        C\ar[r]^\gamma\ar[d]_{i_y} & B_{hg}\ar[d]^{i_x}\\
        C_y\ar[r]_\varphi & B_x
      }
      \tag{6.5.1.1}\label{I.6.5.1.1}
    \]
    is commutative; we therefore have $\varphi=\gamma_x$, and since $\varphi$ is
    a local homomorphism, ${}^a\gamma(x)=y$; $f=({}^a\gamma,\widetilde{\gamma})$
    is therefore an $S$-morphism of the neighbourhood $D(hg)$ of $x$ into $Y$
    which answers the question.
\end{enumerate}

\begin{corollary}[6.5.2]
\label{I.6.5.2}
Under the hypotheses of (\hyperref[I.6.5.1]{6.5.1, (ii)}), if moreover $X$
is of finite type over $S$, we may suppose the morphism $f$ to be of finite
type.
\end{corollary}

This follows from (\hyperref[I.6.3.6]{6.3.6}).

\begin{corollary}[6.5.3]
\label{I.6.5.3}
Suppose the hypotheses of
(\hyperref[I.6.5.1]{6.5.1, (ii)}) are satisfied, and suppose moreover that
$Y$ is integral and $\varphi$ an injective homomorphism. Then we may suppose
that $f=({}^a\gamma,\widetilde{\gamma})$, where $\gamma$ is injective.
\end{corollary}

Indeed, we may suppose $C$ integral
(\hyperref[I.5.1.4]{5.1.4}), hence $i_y$ injective; it then follows from the
diagram (\hyperref[I.6.5.1.1]{6.5.1.1}) that $\gamma$ is injective.

\begin{proposition}[6.5.4]
\label{I.6.5.4}
Let $f=(\psi,\theta):X\to Y$ be a morphism of finite type, $x$ a point of
$X$, $y=\psi(x)$.
\begin{enumerate}
  \item[(i)] For $f$ to be a local immersion at the point $x$
    (\hyperref[I.4.5.1]{4.5.1}), it is necessary and sufficient that
    $\theta_x^\sharp:\mathscr{O}_y\to\mathscr{O}_x$ be surjective.
  \item[(ii)] Suppose moreover that $Y$ is locally Noetherian. For $f$ to be
    a local isomorphism at the point $x$
    (\hyperref[I.4.5.2]{4.5.2}), it is necessary and sufficient that
    $\theta_x^\sharp$ be an isomorphism.
\end{enumerate}
\end{proposition}

\begin{enumerate}
  \item[(ii)] By virtue of (\hyperref[I.6.5.1]{6.5.1}), there then exist an
    open neighbourhood $V$ of $y$ and a morphism $g:V\to X$ such that
    $g\circ f$ (resp. $f\circ g$) is defined and coincides with the identity in
    a neighbourhood of $x$ (resp. $y$), from which one easily deduces that $f$
    is a local isomorphism.
  \item[(i)] The question being local on $X$ and $Y$, we may suppose $X$ and
    $Y$ affine, with respective rings $A$, $B$; we have
    $f=({}^a\varphi,\widetilde{\varphi})$, where $\varphi$ is a ring
    homomorphism $B\to A$ which makes $A$ a $B$-algebra of finite type; we have
    $\varphi^{-1}(\mathfrak{j}_x)=\mathfrak{j}_y$ and the homomorphism
    $\varphi_x:B_y\to A_x$ deduced from $\varphi$ is surjective. Let $(t_i)$
    ($1\leq i\leq n$) be a system of generators of the $B$-algebra $A$;
    the hypothesis on $\varphi_x$ implies that there exist elements $b_i\in B$
    and an element $c\in B-\mathfrak{j}_y$ such that, in the ring of fractions
    $A_x$, we have $t_i/1=\varphi(b_i)/\varphi(c)$ for $1\leq i\leq n$.
    Consequently (\hyperref[I.1.3.3]{1.3.3}), there exists
    $a\in A-\mathfrak{j}_x$ such that, if we set $g=a\varphi(c)$, we also have
    $t_i/1=a\varphi(b_i)/g$ \emph{in the ring of fractions $A_g$}. This being
    so, there exists by hypothesis a polynomial $Q(X_1,\ldots,X_n)$ with
    coefficients in the ring $\varphi(B)$,
\oldpage[I]{152}
    such that $a=Q(t_1,\ldots,t_n)$; set
    \[
      Q(X_1/T,\ldots,X_n/T)
      =P(X_1,\ldots,X_n,T)/T^m,
    \]
    where $P$ is homogeneous of degree $m$. In the ring $A_g$, we have
    \[
      a/1
      =a^mP(\varphi(b_1),\ldots,\varphi(b_n),\varphi(c))/g^m
      =a^m\varphi(d)/g^m
    \]
    where $d\in B$. Since, in $A_g$,
    $g/1=(a/1)(\varphi(c)/1)$ is invertible by definition, so are $a/1$ and
    $\varphi(c)/1$, and we can therefore write
    \[
      a/1=(\varphi(d)/1)(\varphi(c)/1)^{-m}.
    \]
    We conclude that $\varphi(d)/1$ is also invertible in $A_g$. Then set
    $h=cd$; since $\varphi(h)/1$ is invertible in $A_g$,
    the composite homomorphism
    \[
      B\xrightarrow{\varphi}A\longrightarrow A_g
    \]
    factors as
    \[
      B\longrightarrow B_h\xrightarrow{\gamma}A_g
    \]
    (\hyperref[0.1.2.4]{0, 1.2.4}). Let us show that $\gamma$ is
    \emph{surjective}; it suffices to verify that the image of $B_h$ in $A_g$
    contains the $t_i/1$ and $(g/1)^{-1}$. Now, we have
    \[
      (g/1)^{-1}
      =(\varphi(c)/1)^{m-1}(\varphi(d)/1)^{-1}
      =\gamma(c^m/h),
    \]
    and
    \[
      a/1=\gamma(d^{m+1}/h^m),
    \]
    hence
    \[
      (a\varphi(b_i))/1=\gamma(b_i d^{m+1}/h^m),
    \]
    and since
    $t_i/1=(a\varphi(b_i)/1)(g/1)^{-1}$, our assertion is proved. The
    choice of $h$ implies that $\psi(D(g))\subset D(h)$, and the restriction of
    $f$ to $D(g)$ is equal to
    $({}^a\gamma,\widetilde{\gamma})$; since $\gamma$ is surjective, this
    restriction is a closed immersion of $D(g)$ into $D(h)$
    (\hyperref[I.4.2.3]{4.2.3}).
\end{enumerate}

\begin{corollary}[6.5.5]
\label{I.6.5.5}
Let $f=(\psi,\theta):X\to Y$ be a morphism of finite type. Suppose $X$
irreducible, denote by $x$ its generic point, and set $y=\psi(x)$.
\begin{enumerate}
  \item[(i)] For $f$ to be a local immersion at a point of $X$, it is
    necessary and sufficient that
    $\theta_x^\sharp:\mathscr{O}_y\to\mathscr{O}_x$ be surjective.
  \item[(ii)] Suppose moreover that $Y$ is irreducible and locally Noetherian.
    For $f$ to be a local isomorphism at a point of $X$, it is necessary and
    sufficient that $y$ be the generic point of $Y$ (or, what amounts to the
    same (\hyperref[0.2.1.4]{0, 2.1.4}), that $f$ be a
    \emph{dominant} morphism) and that $\theta_x^\sharp$ be an isomorphism (in
    other words, that $f$ be \emph{birational}
    (\hyperref[I.2.2.9]{2.2.9})).
\end{enumerate}
\end{corollary}

It is clear that (i) follows from
(\hyperref[I.6.5.4]{6.5.4, (i)}), taking into account the fact that every
non-empty open subset of $X$ contains $x$; similarly, (ii) follows from
(\hyperref[I.6.5.4]{6.5.4, (ii)}).

\subsection{Quasi-compact morphisms and morphisms locally of finite type}
\label{subsection:I.6.6}

\begin{definition}[6.6.1]
\label{I.6.6.1}
We say that a morphism $f:X\to Y$ is \emph{quasi-compact} if the inverse
image under $f$ of every quasi-compact open subset of $Y$ is quasi-compact.
\end{definition}

Let $\mathfrak{B}$ be a basis of the topology of $Y$ consisting of
quasi-compact open subsets (for example of affine open subsets); for $f$ to be
quasi-compact, it is necessary and sufficient that the inverse image under $f$
of every set of $\mathfrak{B}$ be quasi-compact (or, what amounts to the same,
a \emph{finite} union of affine open subsets), since every quasi-compact open
subset of $Y$ is a finite union of sets of $\mathfrak{B}$. For example, if $X$
is \emph{quasi-compact} and $Y$ \emph{affine}, then \emph{every} morphism
$f:X\to Y$ is quasi-compact: indeed, $X$ is a finite union of affine open
sets $U_i$, and for every affine open subset $V$ of $Y$,
$U_i\cap f^{-1}(V)$ is affine
(\hyperref[I.5.5.10]{5.5.10}), hence quasi-compact.

If $f:X\to Y$ is a quasi-compact morphism, it is clear that for every
open subset $V$ of $Y$, the restriction of $f$ to $f^{-1}(V)$ is a
quasi-compact morphism $f^{-1}(V)\to V$. Conversely, if $(U_\alpha)$ is an
open covering of $Y$ and $f:X\to Y$ a morphism such that the restrictions
$f^{-1}(U_\alpha)\to U_\alpha$ are quasi-compact, then $f$ is
quasi-compact.

\begin{definition}[6.6.2]
\label{I.6.6.2}
We say that a morphism $f:X\to Y$ is \emph{locally of finite type} if, for
every $x\in X$, there exist an open neighbourhood $U$ of $x$ and an open
neighbourhood $V\supset f(U)$ of $f(x)$ such that the
\oldpage[I]{153}
restriction of $f$ to $U$ is a morphism of finite type from $U$ to $V$. We
then also say that $X$ is a prescheme locally of finite type over $Y$, or a
$Y$-prescheme locally of finite type.
\end{definition}

It follows immediately from (\hyperref[I.6.3.2]{6.3.2}) that if $f$ is
locally of finite type, then, for every open subset $W$ of $Y$, the
restriction of $f$ to $f^{-1}(W)$ is a morphism $f^{-1}(W)\to W$ which is
locally of finite type.

If $Y$ is locally Noetherian and $X$ is locally of finite type over
$Y$, then $X$ is locally Noetherian by virtue of
(\hyperref[I.6.3.7]{6.3.7}).

\begin{proposition}[6.6.3]
\label{I.6.6.3}
A morphism $f:X\to Y$ is of finite type if and only if it is quasi-compact and
locally of finite type.
\end{proposition}

The necessity of the conditions is immediate, in view of
(\hyperref[I.6.3.1]{6.3.1}) and the remark following
(\hyperref[I.6.6.1]{6.6.1}). Conversely, suppose these conditions are
satisfied and let $U$ be an affine open subset of $Y$, with ring $A$; for every
$x\in f^{-1}(U)$, there are by hypothesis a neighbourhood
$V(x)\subset f^{-1}(U)$ of $x$ and a neighbourhood $W(x)\subset U$ of
$y=f(x)$, containing $f(V(x))$ and such that the restriction of $f$ to $V(x)$
is a morphism $V(x)\to W(x)$ which is of finite type. Replacing $W(x)$ by a
neighbourhood $W_1(x)\subset W(x)$ of $y$ of the form $D(g)$ (with $g\in A$),
and $V(x)$ by $V(x)\cap f^{-1}(W_1(x))$, we may suppose that $W(x)$ is of the
form $D(g)$, hence of finite type over $U$ (since its ring can be written
$A[1/g]$); consequently $V(x)$ is of finite type over $U$. Moreover,
$f^{-1}(U)$ is quasi-compact by hypothesis, hence a union of finitely many
open subsets $V(x_i)$, which completes the proof.

\begin{proposition}[6.6.4]
\label{I.6.6.4}
\begin{enumerate}
  \item[(i)] An immersion $X\to Y$ is quasi-compact if it is closed, or if the
    underlying space of $Y$ is locally Noetherian, or if the underlying space
    of $X$ is Noetherian.
  \item[(ii)] The composite of two quasi-compact morphisms is quasi-compact.
  \item[(iii)] If $f:X\to Y$ is a quasi-compact $S$-morphism, then so is
    $f_{(S')}:X_{(S')}\to Y_{(S')}$ for every extension $g:S'\to S$
    of the base prescheme.
  \item[(iv)] If $f:X\to X'$ and $g:Y\to Y'$ are two quasi-compact
    $S$-morphisms, then $f\times_S g$ is quasi-compact.
  \item[(v)] If the composite of two morphisms $f:X\to Y$, $g:Y\to Z$ is
    quasi-compact and if $g$ is separated or the underlying space of $X$ is
    locally Noetherian, then $f$ is quasi-compact.
  \item[(vi)] For a morphism $f$ to be quasi-compact, it is necessary and
    sufficient that $f_{\mathrm{red}}$ be so.
\end{enumerate}
\end{proposition}

Note that (vi) is evident, since the property of being quasi-compact for a
morphism depends only on the corresponding continuous map of the underlying
spaces. Let us likewise prove the part of (v) corresponding to the case where
the underlying space $X$ is assumed to be locally Noetherian. Set
$h=g\circ f$, and let $U$ be a quasi-compact open subset of $Y$; $g(U)$ is
quasi-compact (not necessarily open) in $Z$, hence contained in a finite union
of quasi-compact open subsets $V_j$
(\hyperref[I.2.1.3]{2.1.3}), and $f^{-1}(U)$ is consequently contained in
the union of the $h^{-1}(V_j)$, which are quasi-compact subspaces of $X$,
hence Noetherian subspaces. We conclude
(\hyperref[0.2.2.3]{0, 2.2.3}) that $f^{-1}(U)$ is a Noetherian space,
and \emph{a fortiori} quasi-compact.

To prove the other assertions, it suffices to prove (i), (ii) and (iii)
(\hyperref[I.5.5.12]{5.5.12}). Now, (ii) is evident, and (i) follows from
(\hyperref[I.6.3.5]{6.3.5}) when the space $Y$ is locally Noetherian
or the space $X$ Noetherian, and is evident for a closed immersion. To
establish (iii),
\oldpage[I]{154}
we may restrict to the case where $S=Y$ (\hyperref[I.3.3.11]{3.3.11}); set
$f'=f_{(S')}$, and let $U'$ be a quasi-compact open subset of $S'$. For every
$s'\in U'$, let $T$ be an affine open neighbourhood of $g(s')$ in $S$, and let
$W$ be an affine open neighbourhood of $s'$ contained in
$U'\cap g^{-1}(T)$; it will suffice to show that $f'^{-1}(W)$ is
quasi-compact; in other words, we may reduce to proving that, when $S$
and $S'$ are affine, the underlying space of $X\times_S S'$ is
quasi-compact. But since $X$ is then, by hypothesis, a finite union of affine
open subsets $V_j$, $X\times_S S'$ is the union of the underlying spaces of the
affine schemes $V_j\times_S S'$
(\hyperref[I.3.2.2]{3.2.2} and \hyperref[I.3.2.7]{3.2.7}), which
completes the proof of the proposition.

We note also that if $X=X'\amalg X''$ is the sum of two preschemes, a
morphism $f:X\to Y$ is quasi-compact if and only if its restrictions to
$X'$ and $X''$ are.

\begin{proposition}[6.6.5]
\label{I.6.6.5}
Let $f:X\to Y$ be a quasi-compact morphism. For $f$ to be dominant, it is
necessary and sufficient that for every generic point $y$ of an irreducible
component of $Y$, $f^{-1}(y)$ contain the generic point of an irreducible
component of $X$.
\end{proposition}

It is immediate that the condition is sufficient (without supposing $f$
quasi-compact). To see that it is necessary, consider an affine open
neighbourhood $U$ of $y$; $f^{-1}(U)$ is quasi-compact, hence a
\emph{finite} union of affine open subsets $V_i$, and the hypothesis that $f$ is
dominant implies that $y$ belongs to the closure in $U$ of one of the
$f(V_i)$. We may evidently suppose $X$ and $Y$ reduced; since the closure in
$X$ of an irreducible component of $V_i$ is an irreducible component of $X$
(\hyperref[0.2.1.6]{0, 2.1.6}), we may replace $X$ by $V_i$, and $Y$ by
the reduced closed sub-prescheme of $U$ having $\overline{f(V_i)}\cap U$ as
underlying space (\hyperref[I.5.2.1]{5.2.1}), and we are thus reduced to
proving the proposition when $X=\operatorname{Spec}(A)$ and
$Y=\operatorname{Spec}(B)$ are affine and reduced. Since $f$ is dominant,
$B$ is then a subring of $A$ (\hyperref[I.1.2.7]{1.2.7}), and the
proposition then follows from the fact that every minimal prime ideal of $B$
is the intersection of $B$ with a minimal prime ideal of $A$
(\hyperref[0.1.5.8]{0, 1.5.8}).

\begin{proposition}[6.6.6]
\label{I.6.6.6}
\begin{enumerate}
  \item[(i)] Every local immersion is locally of finite type.
  \item[(ii)] If two morphisms $f:X\to Y$, $g:Y\to Z$ are locally of finite
    type, then so is $g\circ f$.
  \item[(iii)] If $f:X\to Y$ is an $S$-morphism locally of finite type, then
    $f_{(S')}:X_{(S')}\to Y_{(S')}$ is locally of finite type for every
    extension $S'\to S$ of the base prescheme.
  \item[(iv)] If $f:X\to X'$ and $g:Y\to Y'$ are two $S$-morphisms
    locally of finite type, then $f\times_S g$ is locally of finite type.
  \item[(v)] If the composite $g\circ f$ of two morphisms is locally of
    finite type, then $f$ is locally of finite type.
  \item[(vi)] If a morphism $f$ is locally of finite type, then so is
    $f_{\mathrm{red}}$.
\end{enumerate}
\end{proposition}

By virtue of (\hyperref[I.5.5.12]{5.5.12}), it suffices to prove (i), (ii)
and (iii). If $j:X\to Y$ is a local immersion, then for every $x\in X$ there
are an open neighbourhood $V$ of $j(x)$ in $Y$ and an open neighbourhood $U$ of
$x$ in $X$ such that the restriction of $j$ to $U$ is a closed immersion
$U\to V$ (\hyperref[I.4.5.1]{4.5.1}), hence this restriction is of finite
type. To establish (ii), consider a point $x\in X$; there are by hypothesis an
open neighbourhood $W$ of $g(f(x))$ and an open neighbourhood $V$ of $f(x)$
such that $g(V)\subset W$ and $V$ is of finite type over $W$; moreover,
$f^{-1}(V)$ is locally of finite type over $V$
(\hyperref[I.6.6.2]{6.6.2}), hence there is an open neighbourhood $U$ of
$x$ which
\oldpage[I]{155}
is contained in $f^{-1}(V)$ and of finite type over $V$; consequently we have
$g(f(U))\subset W$ and $U$ is of finite type over $W$
(\hyperref[I.6.3.4]{6.3.4, (ii)}). Finally, to prove (iii), we may
restrict to the case where $Y=S$ (\hyperref[I.3.3.11]{3.3.11}); for every
$x'\in X'=X_{(S')}$, let $x$ be the image of $x'$ in $X$, $s$ the image of $x$
in $S$, $T$ an open neighbourhood of $s$, $T'$ its inverse image in $S'$,
$U$ an open neighbourhood of $x$ whose image is contained in $T$ and which is
of finite type over $T$; then $U\times_S T'=U\times_T T'$ is an open
neighbourhood of $x'$ (\hyperref[I.3.2.7]{3.2.7}) which is of finite type
over $T'$ (\hyperref[I.6.3.4]{6.3.4, (iv)}).

\begin{corollary}[6.6.7]
\label{I.6.6.7}
Let $X$, $Y$ be two $S$-preschemes which are locally of finite type over $S$.
If $S$ is locally Noetherian, then $X\times_S Y$ is locally Noetherian.
\end{corollary}

Indeed, $X$, being locally of finite type over $S$, is locally Noetherian, and
$X\times_S Y$ is locally of finite type over $X$, hence is also locally
Noetherian.

\begin{remark}[6.6.8]
\label{I.6.6.8}
Proposition (\hyperref[I.6.3.10]{6.3.10}) and its proof extend immediately
to the case where we suppose only that the morphism $f$ is locally of finite
type. Likewise, Propositions
(\hyperref[I.6.4.2]{6.4.2}) and (\hyperref[I.6.4.9]{6.4.9}) remain
valid when we suppose only that the preschemes $X$, $Y$ appearing in their
statement are locally of finite type over the field $K$.
\end{remark}
